\begin{graphillustration}{Пример разбиения Галлаи–Эдмондса и одно максимальное паросочетание (толстые рёбра). $D$: треугольник и две одиночные вершины ($t=3$), $A=\{a\}$, $C=\{c_1,c_2\}$. $\dfc=t-|A|=2$: свободны $d_5$ и $d_3$. Ребро $ac_1$ (между $A$ и $C$) не входит ни в одно максимальное паросочетание.}{fig:06-structure}
\begin{tikzpicture}[font=\small]
\node[gi vertex] (a) at (0,0) {$a$};
\node[gi vertex] (d1) at (-2.4,1.3) {$d_1$};
\node[gi vertex] (d2) at (-3.4,2.0) {$d_2$};
\node[gi free] (d3) at (-3.4,0.6) {$d_3$};
\node[gi vertex] (d4) at (-2.4,-0.4) {$d_4$};
\node[gi free] (d5) at (-2.4,-1.5) {$d_5$};
\node[gi vertex] (c1) at (2.2,0) {$c_1$};
\node[gi vertex] (c2) at (3.6,0) {$c_2$};
\draw[gi edge] (d1) -- (d3); \draw[gi edge] (a) -- (d1);
\draw[gi match] (d1) -- (d2);
\draw[gi edge] (d2) -- (d3);
\draw[gi match] (a) -- (d4); \draw[gi edge] (a) -- (d5);
\draw[gi edge] (a) -- (c1); \draw[gi match] (c1) -- (c2);
\node[gi group,fit=(d1)(d2)(d3)(d4)(d5),label=above:{$D(G)$}] {};
\node[gi group,fit=(a),label=above:{$A(G)$}] {};
\node[gi group,fit=(c1)(c2),label=above:{$C(G)$}] {};
\end{tikzpicture}
\end{graphillustration}
